% Part: normal-modal-logic
% Chapter: frame-correspondence
% Section: properties-accessibility

\documentclass[../../../include/open-logic-section]{subfiles}

\begin{document}

\olfileid[hi]{nml}{frd}{acc}

\olsection{अभिगम्यता संबंधों के गुण}

अनेक मोडल !!{formula} अभिगम्यता संबंध के सरल, यहाँ तक कि परिचित, गुणों के
अभिलाक्षणिक निकलते हैं। एक दिशा में इसका अर्थ है कि कोई दिया हुआ गुण रखने
वाला प्रत्येक प्रतिरूप किसी संगत !!{formula} (और उसके सभी
प्रतिस्थापन-दृष्टांतों) को सत्य बनाता है। हम अभिगम्यता संबंधों के पाँच
शास्त्रीय प्रकारों और उन सूत्रों के उदाहरणों से आरंभ करते हैं जिनकी सत्यता
वे सुनिश्चित करते हैं।

\begin{thm}\ollabel{thm:soundschemas}
  मान लें $\mModel{M} = \tuple{W, R, V}$ एक प्रतिरूप है। यदि $R$ में
  \olref{tab:five} के बाएँ पक्ष का गुण है, तो दाएँ पक्ष के !!{formula} का
  प्रत्येक दृष्टांत $\mModel{M}$ में सत्य है।
\end{thm}

\begin{table}[t]
    \begin{tabular}{| l || l |}
      \hline
      {\emph{यदि $R$ \dots है}} & {\emph{तो \dots $\mModel{M}$ में सत्य है:}} \\
      \hline \hline
      \emph{सीरियल}: $\forall u \exists v Ruv$ & \hbox to.43\textwidth
           {$\Box p \lif \Diamond p$ \hfill (\Ax{D})} \\
      \hline
      \emph{स्वपरावर्ती}: $\forall w Rww$  
      & \hbox to.43\textwidth
          {$\Box p \lif p$ \hfill (\Ax{T})} \\
      \hline
      \emph{सममित}: &  \hbox to .43\textwidth
           {$p \lif \Box\Diamond p$ \hfill (\Ax{B})} \\
           $\forall u\forall v(Ruv \lif Rvu)$ & \\
      \hline
      \emph{संक्रमणशील}: & \hbox to.43\textwidth
           {$\Box p \lif \Box \Box p$ \hfill (\Ax{4})} \\
      $\forall u \forall v \forall w ((Ruv \land Rvw) \lif Ruw)$ & \\
      \hline 
      \emph{यूक्लिडीय}: & \hbox to .43\textwidth
        {$\Diamond p \lif \Box\Diamond p$ \hfill (\Ax{5})} \\
      $\forall w \forall u \forall v ((Rwu \land Rwv) \lif Ruv)$ &\\
      \hline
    \end{tabular}
    \caption{अनुरूपता के पाँच तथ्य।}
    \ollabel{tab:five}
  \end{table} 


\begin{proof}
  \Ax{B} का मामला यह है: स्कीमा की किसी प्रतिरूप में सत्यता दिखाने के लिए
  हमें दिखाना है कि उसके सभी दृष्टांत प्रतिरूप के सभी जगतों पर सत्य हैं।
  अतः $!A \lif \Box\Diamond !A$ को \Ax{B} का कोई दृष्टांत और $w \in W$
  को स्वेच्छ जगत मानें। $\Box \Diamond !A$ की $w$ पर सत्यता दिखाने के लिए
  मान लें कि पूर्ववर्ती $!A$, $w$ पर सत्य है। हमें दिखाना है कि
  $\Diamond !A$, $w$ से अभिगम्य सभी $w'$ पर सत्य है। अब किसी भी $Rww'$ वाले
  $w'$ के लिए सममिति की परिकल्पना से $Rw'w$ भी (देखें
  \olref{fig:Bsymm})। चूँकि $\mSat{M}{!A}[w]$, इसलिए
  $\mSat{M}{\Diamond !A}[w']$। चूँकि $w'$ कोई स्वेच्छ $Rww'$ वाला जगत था,
  हमारे पास $\mSat{M}{\Box\Diamond!A}[w]$ है।

  अन्य मामले अभ्यास के लिए छोड़े गए हैं।
\end{proof}

\begin{prob}
  \olref[nml][frd][acc]{thm:soundschemas} का प्रमाण पूरा कीजिए।
\end{prob}

\begin{figure}
  \begin{center}
    \begin{tikzpicture}[modal]
      \node[world] (w1) [label={below:$\mSat{{}}{\formula{A}}$\\
          $\mSat{{}}{\Box\Diamond \formula{A}}$}] {$w$}; 
      \node[world] (w2) [label=below:{$\mSat{{}}{\Diamond \formula{A}}$},
        right=of w1] {$w'$}; 
      \draw[->, bend left] (w1) to (w2); 
      \draw[->, bend left] (w2) to (w1); 
    \end{tikzpicture}
  \end{center}
\caption{सममिति से तर्क।}
\ollabel{fig:Bsymm}
\end{figure}

ध्यान दें कि \olref{thm:soundschemas} के विलोम निहितार्थ मान्य नहीं हैं:
यह सत्य नहीं कि यदि कोई प्रतिरूप किसी स्कीमा को सत्य बनाता है, तो उस प्रतिरूप
के अभिगम्यता संबंध में संगत गुण होता है। \Ax{T} और स्वपरावर्ती प्रतिरूपों के
मामले में ऐसे प्रतिरूप का उदाहरण देना आसान है जिसमें स्वयं \Ax{T} विफल हो:
$W = \{w\}$ और $V(p) = \emptyset$ लें। तब $R$ स्वपरावर्ती नहीं है, पर
$\mSat{M}{\Box p}[w]$ और $\mSat/{M}{p}[w]$। किंतु यहाँ \Ax{T} का केवल एक
दृष्टांत $\mModel{M}$ में विफल है; अन्य दृष्टांत, जैसे
$\Box \lnot p \lif \lnot p$, सत्य हैं। ऐसे उदाहरण देना कठिन है जिनमें
\Ax{T} का \emph{प्रत्येक प्रतिस्थापन-दृष्टांत} $\mModel{M}$ में सत्य हो और
$\mModel{M}$ स्वपरावर्ती न हो। किंतु ऐसे प्रतिरूप भी हैं:

\begin{prop}\ollabel{prop:reflexive}
  मान लें $\mModel{M} = \tuple{W, R, V}$ ऐसा प्रतिरूप है कि
  $W = \{u, v\}$, जहाँ जगत $u$ और $v$, $R$ से जुड़े हैं: अर्थात् $Ruv$ और
  $Rvu$ दोनों। मान लें कि प्रत्येक $p$ के लिए
  $u \in V(p) \Leftrightarrow v \in V(p)$। तब:
  \begin{enumerate}
  \item प्रत्येक $!A$ के लिए $\mSat{M}{!A}[u]$ तभी जब
    $\mSat{M}{!A}[v]$ ($!A$ पर आगमन का प्रयोग करें)।
  \item \Ax{T} का प्रत्येक दृष्टांत $\mModel{M}$ में सत्य है।
  \end{enumerate}
  चूँकि $\mModel{M}$ स्वपरावर्ती नहीं है (वास्तव में वह
  \emph{अस्वपरावर्ती} है), \Ax{T} के मामले में
  \olref{thm:soundschemas} का विलोम विफल होता है
  (\olref{thm:soundschemas} में उल्लिखित कुछ---पर सभी नहीं---अन्य स्कीमाओं
  के लिए भी ऐसे तर्क दिए जा सकते हैं)।
\end{prop}

\begin{prob}
  \olref[nml][frd][acc]{prop:reflexive} के दावे सिद्ध कीजिए।
\end{prob}

यद्यपि हम पाँच शास्त्रीय !!{formula}ों \Ax{D}, \Ax{T}, \Ax{B}, \Ax{4} और
\Ax{5} पर ध्यान देंगे, फिर भी \olref{tab:anotherfive} में अभिगम्यता संबंधों
के कुछ और गुण दर्ज हैं। यदि प्रत्येक जगत से अधिकतम एक जगत अभिगम्य हो, तो
अभिगम्यता संबंध~$R$ आंशिक फलनीय है। यदि प्रत्येक जगत से ठीक एक जगत अभिगम्य
हो, तो हम उसे फलनीय कहते हैं। (अतः फलनीय संबंध ठीक वे हैं जो सीरियल और
आंशिक फलनीय दोनों हों।) इन्हें ``फलनीय'' इसलिए कहते हैं कि अभिगम्यता संबंध
किसी (आंशिक) फलन की तरह काम करता है। कोई संबंध दुर्बलतः सघन है यदि जब भी
$Ruv$, तब $u$ और~$v$ के ``बीच'' कोई~$w$ हो। इसलिए दुर्बलतः सघन संबंध एक
अर्थ में संक्रमणशील संबंधों के विपरीत हैं: संक्रमणशील संबंध में जब भी आप
$w$ से घूमकर $u$ से~$v$ तक पहुँच सकते हैं, तब $u$ से $v$ तक सीधे भी पहुँच
सकते हैं; दुर्बलतः सघन संबंध में जब भी आप $u$ से~$v$ तक सीधे पहुँच सकते
हैं, तब किसी~$w$ से घूमकर भी पहुँच सकते हैं। कोई संबंध दुर्बलतः दिष्ट है
यदि जब भी किसी जगत~$w$ से जगतों $u$ और~$v$ तक पहुँचा जा सके, तब $u$ और
$v$ दोनों से किसी एक जगत~$t$ तक पहुँचा जा सके---इसे कभी-कभी ``हीरक गुण''
या ``संगमन'' कहते हैं।

\begin{table}[t]
    \begin{tabular}{| p{.50\textwidth} || p{.43\textwidth} |}
      \hline
      {\emph{यदि $R$ \dots है}} & {\emph{तो \dots $\mModel{M}$ में सत्य है:}} \\
      \hline \hline
      \emph{आंशिक फलनीय}: & 
      \multirow{2}{*}{$\Diamond p \lif \Box p$} \\
      $ \forall w \forall u \forall v ((Rwu \land Rwv) \lif u =v )$ & \\
      \hline
      \multirow{2}{*}{\emph{फलनीय}: $\forall w \exists v \forall u(Rwu \liff \eq[u][v]) $} 
      & \multirow{2}{*}{$\Diamond p \liff \Box p$} \\
      & \\
      \hline
      \emph{दुर्बलतः सघन}: &  \multirow{2}{*}{$\Box \Box p \lif
        \Box p$} \\
      $\forall u\forall v(Ruv \lif \exists w(Ruw \land Rwv))$ & \\
      \hline
      \emph{दुर्बलतः संबद्ध}: &
      \multirow{3}{*}{\hbox to.43\textwidth{$
        \begin{array}{@{}l@{}}
          \Box ((p \land \Box p) \lif q) \lor {} \\
          \qquad \Box ((q \land \Box q) \lif p)
        \end{array}$ \hfill (\Ax{L})}
      } \\
      $\forall w \forall u \forall v ((Rwu \land Rwv) \lif {}$ &\\
      \qquad $(Ruv \lor u=v \lor Rvu))$ & \\
      \hline 
      \emph{दुर्बलतः दिष्ट}: & \multirow{3}{*}{\hbox to .43\textwidth
        {$\Diamond\Box p \lif \Box\Diamond p$\hfill (\Ax{G})}} \\
      $\forall w \forall u \forall v ((Rwu \land Rwv) \lif {}$ & \\
      \qquad $\exists t (Rut \land Rvt))$ & \\
      \hline
    \end{tabular}
    \caption{अनुरूपता के पाँच और तथ्य।}
    \ollabel{tab:anotherfive}
  \end{table} 

\begin{prob}
  मान लें $\mModel{M} = \tuple{W, R, V}$ एक प्रतिरूप है। दिखाइए कि यदि
  $R$, \olref[nml][frd][acc]{tab:anotherfive} के बाएँ पक्ष के गुणों को
  संतुष्ट करता है, तो दाएँ पक्ष के संगत सूत्र का प्रत्येक दृष्टांत
  $\mModel{M}$ में सत्य है।
\end{prob}



\end{document}
